% ECONOMICS_PROBLEM: {"id": "Q51-192", "title": "Social cost of carbon under deep uncertainty", "jel_code": "Q51", "source_id": "OP-0111"}
% !TeX program = xelatex
\documentclass[11pt,a4paper]{article}
\usepackage[margin=23mm]{geometry}
\usepackage{fontspec}
\setmainfont{Latin Modern Roman}
\usepackage{amsmath,amssymb,mathtools}
\usepackage{xcolor,enumitem,booktabs,longtable,array,xurl,hyperref}
\definecolor{muted}{HTML}{53636C}
\hypersetup{colorlinks=true,urlcolor=blue,linkcolor=blue}
\setlength{\parindent}{0pt}
\setlength{\parskip}{5pt}
\newcommand{\R}{\mathbb R}
\newcommand{\E}{\mathbb E}
\newcommand{\N}{\mathbb N}
\newcommand{\ind}{\mathbf 1}
\newcommand{\Var}{\operatorname{Var}}
\newcommand{\Cov}{\operatorname{Cov}}
\newcommand{\supp}{\operatorname{supp}}
\DeclareMathOperator*{\argmin}{arg\,min}
\DeclareMathOperator*{\argmax}{arg\,max}
\newcommand{\entrymeta}[3]{{\sffamily\small\color{muted}\textbf{\texttt{#1}}\quad|\quad #2\par #3\par}}
\newcommand{\Problemref}[1]{\texttt{#1}}
\newcommand{\JELref}[2]{\texttt{#2} (JEL \texttt{#1})}
\begin{document}
\section*{Social cost of carbon under deep uncertainty}
\textbf{Problem ID:} \texttt{Q51-192}\quad \textbf{JEL:} \texttt{Q51}\par
% BEGIN ECONOMICS STATEMENT
\label{problem:OP-0111}
\label{atlas:prob:C1}

\noindent\textbf{Original identifier:} C1 (atlas item 111). \textbf{Field:} Environmental and climate economics. \textbf{Classification:} editorial research family with established restricted solutions; current open status not certified.

\subsection*{Definitions and assumptions}

Fix emissions year $t$, a greenhouse-gas species, and a policy baseline. For carbon dioxide let $E_t$ be emissions in metric tonnes of $\mathrm{CO}_2$. A model $m$ and probability law $P$ specify climate dynamics, damage and adaptation processes, population and consumption paths $c_{r,s}>0$ for regions $r$ and dates $s\ge t$. With specified positive welfare weights $\omega_r$, discount factors $\beta_{t,s}$ and differentiable increasing utilities $u_r$, define $W_t(m,P)=\mathbb E_P[\sum_{s=t}^{H}\beta_{t,s}\sum_r\omega_r u_r(c_{r,s})]$. The horizon $H$ is finite unless integrability of the infinite sum is established. For a date-$t$ transfer $T$ with a fixed allocation rule and real-currency numeraire, define $\lambda_t(m,P)=\left.\partial W_t(m,P;T)/\partial T\right|_{T=0}>0$, assumed finite. Its dependence on model and law is retained.

\subsection*{Formal research question}

For an extra emissions pulse $\xi$ measured in tonnes of $\mathrm{CO}_2$, define the model-conditional social cost
\[
\operatorname{SCC}_t(m,P)=
-\lambda_t(m,P)^{-1}
\left.\frac{\partial W_t(m,P;E_t+\xi)}{\partial\xi}\right|_{\xi=0},
\]
provided the derivative and marginal numeraire value are finite. Characterize the range of $\operatorname{SCC}_t$ over an empirically admissible structural class $\mathcal M$ and probability classes $\mathcal P(m)$, allowing uncertain damages, discounting, adaptation and tail behavior. Derive sufficient conditions for finite values, defensible set-valued bounds and mitigation decisions with controlled ambiguity regret. Where differentiation under expectation fails, determine directional welfare changes directly rather than reporting an unsupported finite marginal price.

\subsection*{Interpretation and scope}

Units are real currency per tonne of $\mathrm{CO}_2$; values per tonne of carbon differ by the factor $44/12$. Another greenhouse gas requires its own pulse-response calculation; a chosen $\mathrm{CO}_2$-equivalence convention does not automatically equalize economic damages. Specify whether other policies are fixed or reoptimized after the pulse and how adaptation responds. Welfare weights and the location of the compensating transfer affect monetization across unequal regions. A probabilistic parameter ensemble is not a representation of every structural uncertainty. Robust welfare obtained by an infimum over models may be nondifferentiable, and its derivative need not equal the minimum of individual SCCs.

\subsection*{Source audit and status}

[S1] is a verified early integrated-assessment optimization paper; [S2] analyzes catastrophic-tail uncertainty; [S3] gives probabilistic GIVE estimates using specified evidence and discounting. These sources establish substantial existing model-conditional estimates, not a unique preference-free monetary answer. The structural-class and integrability question above is an editorial refinement. Older publications cannot certify that this precise uncertainty problem remains open in October 2026.

\subsection*{References}

\begin{enumerate}[label={[\textnormal{S}\arabic*]},leftmargin=*]

\item William D. Nordhaus (1992). \emph{An Optimal Transition Path for Controlling Greenhouse Gases}. Science 258(5086), 1315–1319. \url{https://doi.org/10.1126/science.258.5086.1315}. \textbf{Locator:} Publisher or institutional abstract and bibliographic record. \textbf{Support and limits:} Early integrated assessment of optimal mitigation under specified functions and preferences.

\item Martin L. Weitzman (2009). \emph{On Modeling and Interpreting the Economics of Catastrophic Climate Change}. Review of Economics and Statistics 91(1), 1–19. \url{https://doi.org/10.1162/rest.91.1.1}. \textbf{Locator:} Publisher or institutional abstract and bibliographic record. \textbf{Support and limits:} Tail uncertainty and welfare-integrability issues; does not imply every climate model has infinite SCC.

\item Kevin Rennert, Frank Errickson, Brian C. Prest, Lisa Rennels, Richard G. Newell, William Pizer, Cora Kingdon, Jordan Wingenroth, Roger Cooke, Bryan Parthum, David Smith, Kevin Cromar, Delavane Diaz, Frances C. Moore, Ulrich K. Müller, Richard J. Plevin, Adrian E. Raftery, Hana Ševčíková, Hannah Sheets, James H. Stock, Tammy Tan, Mark Watson, Tony E. Wong and David Anthoff (2022). \emph{Comprehensive evidence implies a higher social cost of CO2}. Nature 610, 687–692. \url{https://doi.org/10.1038/s41586-022-05224-9}. \textbf{Locator:} Publisher or institutional abstract and bibliographic record. \textbf{Support and limits:} Probabilistic GIVE estimates with specified discounting and damage components; not unique preference-free SCC.

\end{enumerate}

\subsection*{Common conventions: Source mathematical conventions}

Unless an entry states otherwise, agent and firm sets are finite. State and action spaces are measurable, strategies and outcome maps are measurable, probability laws are specified, and expectations exist for the targets under discussion. Infinite-horizon discount factors lie in $(0,1)$ and discounted objectives are finite. Norms, tolerances and complexity input sizes must be specified for computational comparisons. Constants or primitives that appear locally are inputs, not empirical facts asserted by this catalogue.

\textbf{Identification.} A statistical model $m\in\mathcal M$ implies an observable law $P_m$ and target $\theta(m)$. At law $P$, its identified set is
\[
\Theta(P;\mathcal M)=\{\theta(m):m\in\mathcal M,\ P_m=P\}.
\]
Point identification holds when this set is a singleton, or equivalently when $P_{m_1}=P_{m_2}$ implies $\theta(m_1)=\theta(m_2)$ for all compatible models. Sharp bounds mean bounds attained, or approached when the boundary is not attained, by compatible models. Writing this set is a definition, not a solution: the substantive task is to establish informative restrictions and derive or estimate the set. An unrestricted-confounding model may give only trivial bounds.

\textbf{Inference.} Identification, estimability and valid uncertainty quantification are separate questions. Fix $\alpha\in(0,1)$. A confidence procedure $C_n$ over a declared law class $\mathcal P$ has asymptotic uniform coverage at level $1-\alpha$ when
\[
\liminf_{n\to\infty}\inf_{P\in\mathcal P}\Pr_P\{\theta(P)\in C_n\}\geq1-\alpha.
\]
For set-identified targets the coverage event must be defined separately, for example coverage of the full identified set. If an entry uses identification-indexed models instead of uniquely indexed laws, the same coverage convention is applied over those models. Pointwise limits do not establish uniform validity, and asymptotic validity does not guarantee finite-sample precision. Sample sizes, dependence structures and weak-identification sequences are part of each application.

\textbf{Counterfactuals.} $Y_i(a)$ is a potential outcome under a specified intervention, including its timing, implementation and versions. It is not $Y_i$ conditional on voluntarily choosing $a$. A full intervention vector permits interference; shorthand scalar treatment notation does not silently impose no interference. Assignment, selection, attrition, anticipation and measurement must be specified. A claim about an instrument's local effect is not automatically a population policy effect.

\textbf{Economic equilibrium.} A counterfactual model includes best responses, constraints, feasibility, market clearing, state transitions and expectations consistent with the stated information. When several equilibria exist, either a declared selector is an input or the target is set-valued. Comparisons across policy regimes must use the stated selection convention. Reduced-form relationships are not presumed invariant after an equilibrium-changing intervention.

\textbf{Optimization and welfare.} Optimization is over the stated feasible set. If attainment is not established, $\sup$ or $\inf$ and $\varepsilon$-optimal choices replace an asserted maximizer or minimizer. For example, with value $V=\sup_{a\in\mathcal A}W(a)<\infty$, an $\varepsilon$-optimal set is $\{a\in\mathcal A:W(a)\geq V-\varepsilon\}$ for $\varepsilon>0$. Benefits, costs, risk and health or environmental outcomes can be aggregated only using declared units and conversion weights. Interpersonal utility weights, the treatment of fiscal transfers and foreign profits, discounting, and the behavioral welfare criterion are explicit inputs. A comparative-static derivative additionally requires existence and appropriate differentiability of the selected counterfactual.

\textbf{Sources.} Local labels [S1], [S2], and so forth restart in every question. Locators identify the relevant abstract, model, section or result at the level actually checked; a bibliographic or abstract-level verification is not represented as inspection of an entire proof. Working papers are distinguished from later publications and changing versions. Source summaries are paraphrased. All URL checks and editorial formulations refer to the audit date stated above.
% END ECONOMICS STATEMENT
\end{document}
